%% kapa: exam=1 round=30 year=2019 period=2 form=A subject=appraisal_law subject_ko=감정평가관계법규 seq=13 number=13
%% source: Q-Net 감정평가사 기출문제 (cst00309) artlSeq=5205205 file=extracted/1008/5205205/2019년 제30회 감정평가사 1차시험 문제지/2019년 제30회 감정평가사 1차시험 문제지/2019년도 제30회 감정평가사 1차 2교시 A형.hwp
%% conversion: kordoc 4.24.51 -> tools/gen1.py (auto); review=proofread-sonnet; flags=none
\begin{question}{13}
국토의 계획 및 이용에 관한 법령상 개발행위에 따른 공공시설의 귀속에 관한 설명으로 옳은 것은?
\begin{choices}
  \item 개발행위허가를 받은 자가 행정청이 아닌 경우 개발행위로 용도가 폐지되는 공공시설은 새로 설치한 공공시설의 설치비용에 상당하는 범위에서 개발행위허가를 받은 자에게 무상으로 양도할 수 있다.
  \item 개발행위허가를 받은 자가 행정청이 아닌 경우 개발행위허가를 받은 자가 새로 설치한 공공시설은 그 시설을 관리할 관리청에 유상으로 귀속된다.
  \item 개발행위허가를 받은 자가 행정청인 경우 개발행위허가를 받은 자가 새로 설치한 공공시설은 개발행위허가를 받은 행정청에 귀속된다.
  \item 군수는 공공시설인 도로의 귀속에 관한 사항이 포함된 개발행위허가를 하려면 미리 기획재정부장관의 의견을 들어야 한다.
  \item 개발행위허가를 받은 자가 행정청인 경우 개발행위허가를 받은 자가 준공검사를 마쳤다면 해당 시설의 관리청에 공공시설의 종류를 통지할 필요가 없다.
\end{choices}
\end{question}
\answer{1}
